% Generated draft: requires mathematical and source review.
\documentclass[11pt]{article}
\usepackage[margin=25mm]{geometry}
\usepackage{fontspec}
\setmainfont{Latin Modern Roman}
\usepackage{amsmath,amssymb,mathtools,unicode-math}
\setmathfont{Latin Modern Math}
\usepackage{enumitem,needspace,xurl,hyperref}
\setlength{\parindent}{0pt}
\setlength{\parskip}{5pt}
\setlength{\emergencystretch}{3em}
\newcommand{\N}{\mathbb{N}}
\newcommand{\Z}{\mathbb{Z}}
\newcommand{\Q}{\mathbb{Q}}
\newcommand{\R}{\mathbb{R}}
\newcommand{\C}{\mathbb{C}}
\newcommand{\F}{\mathbb{F}}
\newcommand{\A}{\mathbb{A}}
\newcommand{\PP}{\mathbb{P}}
\newcommand{\E}{\mathbb{E}}
\newcommand{\eps}{\varepsilon}
\newcommand{\dd}{\,\mathrm{d}}
\begin{document}
\section*{UnsolvedMath ID: 20002007}
\textbf{Source problem code:} catalogue entry
\subsection{Definitions and mathematical statement}
\textbf{Section C.2, workshop proposals moderated and contributed by Margaret Symington.} The source asks for intrinsic tropical varieties and abstract idempotent varieties: appropriate local models, charts and gluing, possibly in ringed-space or non-Archimedean analytic language. It further asks when a complex tropical limit is homeomorphic or homotopy equivalent to a nearby algebraic variety; for a Hodge theory compatible with classical limiting Hodge theory; for a refinement/patchworking theorem giving information about a single classical hypersurface; and for maps, line bundles and vector bundles. These are several foundational questions, not a single settled definition or conjecture.

\textbf{A convention-explicit interpretation of the proposed coordinates.} Take the formal Puiseux field $K=\bigcup_{r\ge1}\C((t^{1/r}))$ and an ideal $I\subset K[z_1^{\pm1},\ldots,z_n^{\pm1}]$, with $X=V(I)\subset(K^*)^n$. For $z=c t^\lambda+\text{higher terms}$, $c\ne0$, define $\operatorname{val}(z)=\lambda\in\Q$ and $\operatorname{ph}(z)=c/|c|\in S^1$. The raw tropical image is $\operatorname{val}(X)\subset\Q^n$; the usual closed tropical support is its closure in $\R^n$. A phase image lies in $\Q^n\times(S^1)^n$, or after closure in $\R^n\times(S^1)^n$. A chosen encoding $(w,\theta)\mapsto(e^{-w_j+i\theta_j})_j$ identifies the latter ambient space with $(\C^*)^n$; this encoding is a convention, not the raw valuation map itself. Real phases are the signs $\{\pm1\}^n$, equivalently $\theta_j\in\{0,\pi\}$, rather than the vacuous statement that numerical arguments are real.

The extraction's $K=\C(t)$ called Puiseux series, ideals ``in $K^*$'', and valuation targets $\R^*$ or $\Q^*$ are not adopted literally: an ideal belongs to a ring, a Puiseux series field differs from the rational-function field, and valuation zero must be allowed. Raw rational-valued images and their closures must also be distinguished.

\textbf{Supplied correction of the torus-fiber assertion.} The workshop says that each complex tropical projection fiber is a torus. The aid instead identifies the fiber over $w$, in the appropriate phase/initial-degeneration model, with the coamoeba or closed coamoeba of $\operatorname{in}_wX$. Here a coamoeba is the argument image in $(S^1)^n$. Closure and lifting conventions must be fixed; general coamoebas need not be tori.

For a hypersurface initial form $c_a z^a+c_b z^b$, put $u=a-b\in\Z^n\setminus\{0\}$ and $g=\gcd(|u_1|,\ldots,|u_n|)$. Its exact phase fiber is
\[
 \{\theta\in(\R/2\pi\Z)^n:
      e^{i\langle u,\theta\rangle}
         =\operatorname{ph}(-c_b/c_a)\}.
\]
It has exactly $g$ connected components, each an $(n-1)$-torus, and is a single torus precisely when $u$ is primitive. For real coefficients, writing $\operatorname{sgn}(-c_b/c_a)=(-1)^\tau$, its real phase subset is
\[
 \{s\in\F_2^n:(u\bmod2)\cdot s=\tau\}.
\]
It is empty if that affine equation is inconsistent; otherwise it has $2^{n-1}$ points when $u\bmod2\ne0$, and $2^n$ points when $u\bmod2=0$ and $\tau=0$. The aid gives $1+x+y$ at the tropical-line vertex as a non-torus coamoeba fiber, so even the binomial union-of-tori description is not universal.

\textbf{Remaining source qualifications.} The $p$-adic discussion proposes valuation plus residue/phase data, informally written $\Q\times\overline{\F}_p^*$, but does not define its angular component, topology or the stated ``closure of multiplicative generators''. No canonical phase choice is inferred. Viewing an ideal over formal Puiseux series as a family $I_t$ over complex numbers additionally requires convergence or algebraic family data. A Hausdorff-limit assertion needs rescaled logarithmic coordinates and, for noncompact spaces, a local Hausdorff or compactification convention. The source's claim that a smooth complex tropical curve is homeomorphic to the small nonzero fibers is retained as a claim needing a precise smoothness/degeneration hypothesis, not an unrestricted theorem. The aid mentions a qualified modern status synthesis for intrinsic spaces, analytification/skeleta, Hodge theory, topology, patchworking and bundles, but supplies no specific references or complete resolutions for those questions here.
\subsection{Short English statement}
The workshop asks for intrinsic tropical foundations, topology and Hodge theory of limits, patchworking, maps and bundles. The supplied aid corrects one concrete assertion: phase fibers are initial-degeneration coamoebas, not always tori. A binomial fiber has as many torus components as the gcd of its exponent difference, while a trinomial already gives a non-torus example.
\subsection{Sources}
\begin{itemize}
\item Local UnsolvedMath registry, permanent ID 20002007, source code catalogue entry; canonical \texttt{aim-geometry-notes.json}, notes[344].
\item AIM workshop report \emph{Amoebas and Tropical Geometry}, 14 January 2004, Section C.2 ``Basic definitions'', moderated and contributed by Margaret Symington: \url{https://aimath.org/WWN/amoebas/amoebas.pdf}. The whiteboard figure is mentioned but not supplied as additional evidence.
\item The source records workshop proposals rather than a settled definition. No separate authoritative references for modern phase tropicalization, initial degenerations, topology or Hodge theory are supplied in this record.
\item Registry research attempt 1, dated 2026-08-14: partial result, rigor ``proved'', novelty confidence low. The machine-generated, non-peer-reviewed aid requires independent verification before citation.
\end{itemize}
\end{document}
